\chapter{Family--defect quiver moment-map gate}

Gauge--family locking gate выделил специальный square
\begin{equation}
 \left(|\Phi|^2-\frac13\Tr X^TX\right)^2,
\end{equation}
но не вывел ratio его трёх monomials. Здесь показывается, что этот square
является central частью canonical moment-map norm трёхузлового quiver.
Одновременно тот же functional штрафует anisotropy locking field и в
unit-momentum defect model даёт устойчивый nonzero pairing saddle.

\section{Frozen--gauged--frozen chain}

Рассмотрим три copies real family triplet:
\begin{equation}
 V_L\xrightarrow{\ X\ }V_G\xrightarrow{\ Y\ }V_R.
 \label{eq:family-moment-map-quiver}
\end{equation}
Левый и правый nodes являются frozen global family frames, а middle node
несёт gauged \(SO(3)_G\). Первый arrow \(X\) является real
bifundamental locking field предыдущей главы. Второй arrow несёт ordinary
charge-two \(B-L\) pairing connector.

После tetrahedral reduction обе frozen copies несут один и тот же
irreducible real \(A_4\)-triplet. Прямой commutant audit всех двенадцати
rotations даёт
\begin{equation}
 \dim_{\mathbb R}\operatorname{End}_{A_4}(\mathbb R^3)=1.
\end{equation}
Следовательно, Schur lemma фиксирует equivariant connector:
\begin{equation}
 Y=\Phi I_3.
 \label{eq:a4-schur-pairing-connector}
\end{equation}
Это не isotropic ansatz, выбранный после vacuum: внутри declared
\(A_4\)-equivariant arrow space других directions нет.

\section{Middle-node moment map}

Пусть \(d\) --- oriented quiver differential с blocks \(X\) и \(Y\).
Middle block graded Hodge commutator равен
\begin{equation}
 \mu_G
 =[d,d^\dagger]\big|_{V_G}
 =XX^T-Y^\dagger Y
 =XX^T-|\Phi|^2I_3.
 \label{eq:middle-quiver-moment-map}
\end{equation}
Поскольку только middle node gauged, его auxiliary/D-term action использует
normalized trace
\begin{equation}
 S_\mu=\tau_3(\mu_G^2),
 \qquad
 \tau_3(M)=\frac13\Tr M.
 \label{eq:normalized-middle-moment-action}
\end{equation}

Incoming-minus-outgoing sign является стандартным real moment map quiver
gauge theory. D-term interpretation и frozen flavor nodes описаны, например,
в \path{arXiv:hep-th/0405230}; norm-square moment-map dynamics --- в
\path{arXiv:0807.4734}. Graded connections и quiver Higgs systems из
equivariant gauge reduction рассмотрены в \path{arXiv:0905.2338}. В
finite spectral geometry Yang--Mills curvature norm действительно становится
Higgs potential; см. \path{arXiv:hep-th/0101217}. Эти источники фиксируют
архитектуру; конкретный \(A_4\)-Schur reduction является project result.

\section{Точная coefficient identity}

Обозначим
\begin{equation}
 G=XX^T,
 \qquad
 \bar g=\tau_3(G),
 \qquad
 G_0=G-\bar gI_3.
\end{equation}
Тогда \(\tau_3(G_0)=0\), и central/traceless orthogonality даёт
\begin{equation}
 \boxed{
 \tau_3(\mu_G^2)
 =\left(|\Phi|^2-\frac13\Tr XX^T\right)^2
 +\frac13\Tr G_0^2.}
 \label{eq:moment-map-central-shape-decomposition}
\end{equation}
Первое слагаемое является ровно искомым norm-locking square, второе ---
положительный penalty за departure от isotropic frame.

Expansion central term имеет coefficient ratio
\begin{equation}
 |\Phi|^4:
 |\Phi|^2\Tr XX^T:
 (\Tr XX^T)^2
 =1:-\frac23:\frac19.
 \label{eq:derived-pairing-ratio}
\end{equation}
Ratio поэтому следует не из настройки portal, а из dimension три и
normalized trace middle node. Identity проверена на 128 random matrices;
максимальный residual меньше \(2.2\times10^{-14}\).

\begin{theorem}[Quiver moment-map coefficient pass]
Frozen--gauged--frozen family quiver с \(A_4\)-equivariant right arrow
однозначно даёт moment map
\(\mu_G=XX^T-|\Phi|^2I_3\). Его normalized curvature norm содержит
pairing square с exact ratio \(1:-2/3:1/9\) и дополнительный positive
shape-locking term. Independent portal coefficients отсутствуют.
\end{theorem}

\section{Vacuum Hessian}

Объединим \(S_\mu\) с уже проверенным bifundamental potential
\begin{equation}
 V_{\rm lock}
 =\frac14\|X^TX-I_3\|_F^2
 +\frac14(\det X-1)^2.
\end{equation}
В zero-gradient vacuum
\begin{equation}
 X_\star=I_3,
 \qquad
 |\Phi_\star|=1
\end{equation}
full eleven-real-dimensional Hessian имеет signature
\begin{equation}
 (7_+,4_0,0_-).
\end{equation}
Три zero modes являются gauge orbit locking field, четвёртая --- phase
pairing field до включения winding/gauge data. Отрицательных modes нет.

\section{Unit-momentum defect test}

Ordinary charge-two field на surviving root branch имеет lowest covariant
momentum \(k^2=1\) в normalized geometry \(L=\pi\). Чтобы не замораживать
frame radius вручную, рассмотрим joint radial slice
\begin{equation}
 X=\rho I_3,
 \qquad
 \Phi=r e^{i\varphi},
\end{equation}
с energy
\begin{equation}
 V_{\rm rad}(\rho,r)
 =V_{\rm lock}(\rho I_3)
 +(\rho^2-r^2)^2+r^2.
 \label{eq:joint-moment-map-defect-radial}
\end{equation}
Eight-start global numerical minimization даёт
\begin{equation}
 \rho_\star=0.7432242844,
 \qquad
 r_\star=0.2288718801,
 \qquad
 V_\star=0.5395181198.
\end{equation}
Лучший normal stationary branch имеет
\begin{equation}
 r=0,
 \qquad
 \rho=0.72442382,
 \qquad
 V_{\rm normal}=0.5408201282.
\end{equation}
Следовательно, condensed branch ниже на
\begin{equation}
 \Delta V=0.0013020084.
\end{equation}
Radial Hessian condensed saddle имеет eigenvalues
\begin{equation}
 0.19471629,
 \qquad
 8.67361589,
\end{equation}
и потому положителен.

\begin{proposition}[Normalized defect-condensate pass]
При canonical unit coefficients common locking-plus-moment-map action
сохраняет stable nonzero pairing condensate даже после включения lowest
unit defect momentum. Condensation является слабой, но energetic ordering
однозначно положителен внутри проверенной normalized model.
\end{proposition}

\section{Граница результата}

Коэффициентный gate закрыт, но physical parent ещё требует actual finite
embedding. Нужно показать, что:
\begin{enumerate}
 \item nodes \(V_L,V_G,V_R\), их frozen/gauged status и arrows возникают
 из одной finite graded algebra и real structure;
 \item right arrow действительно имеет charge \(-2\) по \(B-L\) и
 Schur-restricted form \(\Phi I_3\);
 \item middle-node moment-map norm входит с той же normalization, что
 scalar kinetic term и \(V_{\rm lock}\);
 \item KO6 doubling, BV fields и nonradial fluctuations не меняют
 ordering condensed/normal branches;
 \item полученный background сохраняет exact-one BdG kernel.
\end{enumerate}

\begin{nogocriterion}[Quiver embedding gate]
Нельзя объявлять pairing condensate выведенным из spectral geometry, пока
moment-map quiver предъявлен как external gauge-theory block. Следующий pass
требует построить его как finite Dirac/superconnection complex и получить
normalization~\eqref{eq:normalized-middle-moment-action} из того же trace,
который нормирует kinetic sector.
\end{nogocriterion}

Итак, portal ratio и normalized condensate больше не являются свободными
EFT choices. Остаётся один representation-level вопрос: существует ли этот
трёхузловой moment-map block внутри фактической finite geometry проекта.

Следующий KO6 embedding gate отвечает на этот вопрос частично положительно.
Explicit 18-dimensional chain проходит reality, grading, order-zero и
first-order conditions; последний сам фиксирует \(Y=\Phi I_3\). Однако
ordinary spectral quartic имеет mixed sign, противоположный moment-map
action. Поэтому representation embedding закрыт, а action embedding
переносится в relative/auxiliary curvature class.

Последующий relative/auxiliary audit уточняет terminology. Canonical
mapping-cone norm выбирает только endpoint product, а strict
\(SO(3)\)-adjoint projection симметричной \(\mu_G\) равна нулю. Поэтому
\(S_\mu\) не является обычным \(SO(3)\) D-term. Algebraic pass сохраняется
для self-adjoint auxiliary module \(\operatorname{Sym}_3(\mathbb R)\), но
его degree-two calculus/BV origin остаётся следующим gate.

Машинный ledger сохранён в
\path{s2t_v4_family_defect_quiver_moment_map_gate_results.json}; генератор
--- \path{s2t_v4_family_defect_quiver_moment_map_gate.py}.